%------------------------------------------------------------------------------%
\section{Função Injetora}

%------------------------------------------------------------------------------%
\begin{definicao}{}{def:}(\emph{Função Injetora}):
  Uma função $f$, definida em seu domínio $D$ é \emph{injetora} quando dados
  quaisquer valores reais $x_1$, $x_2 \in D$, se $x_1 \neq x_2$, então, $f(x_1)
  \neq f(x_2)$.
\end{definicao}
%------------------------------------------------------------------------------%

Em outras palavras, se $f$ é uma função injetora então,
\[
  f(x_1)=f(x_2) \Leftrightarrow x_1=x_2.
\]

\emph{Teste da reta horizontal:}
Uma função é injetora se e somente se nenhuma reta
horizontal intercepta seu gráfico mais de uma vez.

%------------------------------------------------------------------------------%
\begin{figure}
  \centering
  \includegraphics[scale=1.4]{injet}
  \caption{Gráfico de uma função que não é injetora,
  uma vez que $f(x_1)=f(x_2)$ para $x_1 \neq x_2$.}
\end{figure}
%------------------------------------------------------------------------------%

%------------------------------------------------------------------------------%
\begin{example}
  Verifique se a função \(f(x) = x^2\) é injetora

  \solution
  $f$ não é injetora pois $f(-1)=f(1)=1$.
\end{example}
%------------------------------------------------------------------------------%

\emph{Observação sobre Contraexemplos:}
Muitas vezes, é mais simples demonstrar que uma função \emph{não} possui
determinada propriedade do que provar que ela possui.
Para mostrar que uma função não é injetora, por exemplo, basta apresentar um
\emph{contraexemplo}\index{Contraexemplo}.

Um contraexemplo é um caso específico que invalida uma afirmação geral.
No exemplo anterior, para mostrar que $f(x)=x^2$ não é injetora,
bastou encontrar dois elementos distintos ($x_1=1$ e $x_2=-1$)
que resultam na mesma imagem ($f(1)=f(-1)=1$).
Se a regra falha para um par, ela deixa de ser válida para todo o domínio.

%------------------------------------------------------------------------------%
\begin{example}
  Verifique se a função \(f(x) = x^3\) é injetora

  \solution
  Suponha que
  \[
    f(x_1)=f(x_2)
  \]
  temos então que
  \begin{align*}
    x_1^3 & = x_2^3 \\
    x_1^3 & = \sqrt[3]{x_2^3} \\
    x_1   & = x_2,
  \end{align*}
  logo $f$ é injetora.
\end{example}
%------------------------------------------------------------------------------%

%------------------------------------------------------------------------------%
\begin{example}
  Verifique se a função \(y= \frac{3}{5x-2}\) é injetora

  \solution
  Suponha que
  \[
    f(x_1)=f(x_2)
  \]
  temos então que
  \begin{align*}
   \frac{3}{5x_1-2} & = \frac{3}{5x_2-2} \\
   3(5x_1-2)        & = 3(5x_2-2) \\
   5x_1-2           & = 5x_2-2\\
   x_1              & = x_2.
  \end{align*}
  Logo, $y$ é injetora.
\end{example}
%------------------------------------------------------------------------------%